\chapter{Endpoints --- Communication as a First-Class Primitive}

An endpoint answers: \textbf{``Who am I communicating with?''} It is
CharlotteOS's primary primitive for calls between protection domains.

\section{Server and connection model}

An \textbf{endpoint} is a bounded server-side IPC receive object. A service creates
one and receives messages on it. A \textbf{connection capability} is a client-side
authority to send to that endpoint. Connections are minted from the endpoint
(typically by the name service) and delegated to clients.

\begin{verbatim}
Service creates endpoint:   endpoint_create(interface="ns", capacity=16)
Name service mints conn:    connection_mint(endpoint, SEND | CALL)
Client receives conn:       name_service.lookup("ns") -> ConnectionCap
\end{verbatim}

Ordinary clients do not mint arbitrary connections. The name service or a policy
service normally delegates them after an access check.

\section{Message classes}

The ABI distinguishes two message classes:

\subsection{Scalar messages}

A scalar message carries an opcode, one 64-bit argument word, and optionally a
reply token. It fits in registers. It is suitable for:

\begin{itemize}
    \item Opcodes and request types.
    \item Capability handles.
    \item Status values and queue indices.
    \item Compact control-plane operations where the payload is a single word.
\end{itemize}

Scalar messages require no memory-object allocation, page-table change, or TLB
invalidation.

\subsection{Memory-bearing messages}

A message can carry one or more \textbf{memory-object attachments} with explicit
transfer modes (see Chapter~6). These are suitable for:

\begin{itemize}
    \item Structured requests and responses.
    \item Strings, names, and configuration structures.
    \item Packet payloads and block I/O buffers.
    \item Larger replies that exceed register width.
\end{itemize}

The kernel validates each attachment: the sender must hold the referenced memory
object with sufficient rights for the declared transfer mode. The receiver gets a
new capability to the memory (or a temporary mapping, for borrows).

\section{Message envelope}

The kernel validates a fixed-width message envelope. It does not understand Rust
enums, \texttt{serde} formats, or protocol-specific types:

\begin{verbatim}
struct MessageHeader {
    interface: u128,       // protocol identity
    version: u32,
    opcode: u32,
    flags: u32,
    inline_len: u16,
    segment_count: u16,
    capability_count: u16,
    reserved: u16,
    user_tag: u64,
}
\end{verbatim}

Constraints on the ABI:

\begin{itemize}
    \item Fixed-width fields --- no \texttt{usize}, no pointer-sized values.
    \item Explicit alignment --- no compiler-dependent layout.
    \item Checked offsets and lengths --- the kernel rejects out-of-bounds.
    \item Explicit protocol versioning --- forward compatibility is negotiated.
    \item Maximum attachment counts --- bounded kernel work per message.
    \item No arbitrary userspace pointers as durable message identity.
\end{itemize}

\section{Call and reply}

The most common IPC pattern is \textbf{call/reply}: a client sends a request and
expects a response.

\subsection{Synchronous-looking futures}

At the Rust level, a client writes:

\begin{verbatim}
let response = network_service.connect(address).await?;
\end{verbatim}

Underneath, the runtime:

\begin{enumerate}
    \item Sends a call message through the connection capability.
    \item The kernel creates a one-shot \textbf{reply token} bound to this call.
    \item The caller task becomes pending (the \texttt{Future} returns
        \texttt{Poll::Pending}).
    \item The caller shard continues running other tasks --- no thread is
        dedicated to this request.
    \item The server receives the message on its endpoint.
    \item The server replies (immediately or later).
    \item Reply arrival wakes the client future, which returns
        \texttt{Poll::Ready(response)}.
\end{enumerate}

\subsection{Deferred replies}

A server may retain a reply token and complete it later --- while continuing to
serve other requests on the same shard. This supports device-driven
operations:

\begin{verbatim}
async fn handle_read(pending: PendingRequest) {
    // Submit a DMA read, retaining the reply token.
    // Continue serving other clients while the read is in flight.
    let buffer = device.read(address, len).await?;
    pending.reply_with_buffer(buffer);
}
\end{verbatim}

The shard's executor polls other tasks while the read completes.

\subsection{Reply token semantics}

A reply token is:

\begin{itemize}
    \item \textbf{Created by the kernel} when a call message is enqueued.
    \item \textbf{Bound to the pending call} --- it completes exactly that
        caller's future.
    \item \textbf{Consumable exactly once} --- a second reply attempt returns
        an error.
    \item \textbf{Invalidated by cancellation} --- if the caller times out or
        aborts, the reply token becomes inert.
    \item \textbf{Invalidated by endpoint teardown} --- if the server dies,
        all outstanding reply tokens are invalidated, and all waiting
        callers receive \texttt{ServiceRestarted}.
    \item \textbf{Optionally delegable} --- a server may forward a reply token
        to another service, allowing a chain of services to contribute to
        a single reply.
\end{itemize}

\section{Connection delegation}

A connection capability can be delegated to another domain with attenuated rights.
This is the mechanism by which the userspace name service returns service-level
connections:

\begin{verbatim}
// Name service:
let full_connection = name_service.lookup("driver.nic.v1")?;
// full_connection has SEND | CALL | DELEGATE rights

// Delegate to a client with reduced rights:
connection_delegate(full_connection, client_domain, SEND | CALL)?;
// client_domain now holds a connection with SEND | CALL only
\end{verbatim}

The kernel mediates. The name service never sees raw endpoint identities; it
holds connection capabilities, which it re-delegates. Each delegation step can
only reduce rights.

\section{Service identity and discovery}

The kernel uses opaque endpoint identities. A userspace \textbf{name service}
maps human-readable names to re-delegable connection capabilities:

\begin{itemize}
    \item \textbf{Registration:} a service registers as \texttt{"driver.nic.v1"}
        with a specific endpoint. The name service stores the endpoint's
        connection.
    \item \textbf{Generation tracking:} each registration bumps a generation
        counter. Stale client connections bound to the previous generation fail
        with \texttt{ServiceRestarted}.
    \item \textbf{Access-key gating:} a registration may specify a non-zero access
        key. Lookups without the matching key are denied --- policy enforced by
        a userspace service, not the kernel.
    \item \textbf{Lookup:} a client asks for \texttt{"driver.nic.v1"} and receives
        a connection capability (or an error). It never receives an address, a
        port number, or a raw handle.
\end{itemize}

The name service runs in userspace. The kernel provides the endpoint and
connection primitives; userspace provides the naming policy.

\section{Two kernel data paths, not one}

A central distinction is that endpoint IPC is for \textbf{cross-domain
service calls}. Completion queues are for \textbf{kernel and device operations}.
They are not the same mechanism.

\begin{center}
\begin{tabular}{|l|l|}
\hline
\textbf{Endpoint IPC path} & \textbf{Completion queue path} \\
\hline
Application $\rightarrow$ network service & Timer expiry \\
Network service $\rightarrow$ NIC driver & DMA completion \\
Application $\rightarrow$ filesystem service & IRQ-derived device completion \\
Service $\rightarrow$ name service & Waiting for thread/process exit \\
\hline
\end{tabular}
\end{center}

A userspace driver may compose both: receive an endpoint message (data path), then
submit a hardware operation and await a CQ entry (control path). The endpoint
message and the hardware completion are related \emph{in the service's logic}, but
they are different kernel abstractions with different semantics.

\section{Why endpoints matter for critical software}

\begin{enumerate}
    \item \textbf{Communication carries authority.} A client that
        holds a connection capability was authorized (at grant time, by the name
        service or policy service) to communicate with the target. Every message
        is authorized by possession of a kernel-mediated capability rather than
        by an address or UID alone.

    \item \textbf{Scalar and memory messages are distinct.} The kernel does not
        copy arbitrary data between address spaces by default. Small control
        messages stay small. Bulk data moves as memory objects (Chapter~6) with
        explicit ownership semantics.

    \item \textbf{Deferred replies enable efficient concurrency.} A server is never
        forced to block waiting for one slow client. It retains reply tokens for
        in-flight work and continues serving other requests. The shard's executor
        multiplexes naturally.

    \item \textbf{Restart errors are defined.} Stale connections return typed
        errors such as \texttt{EndpointClosed} and \texttt{ServiceRestarted}.
        A replacement cannot inherit the old instance's reply tokens.

    \item \textbf{Name service is policy, not mechanism.} The kernel provides
        endpoints and connections. Userspace decides who can register what name,
        who can look up what service, and what access keys are required. The
        kernel remains small; policy is replaceable without a kernel rebuild.
\end{enumerate}

\endinput
